\section{Intervals with absolute index prime to $p$}\label{sec:prime-to-p}
This section proves the uniform prime-to-$p$ character bound by recovering canonical lines and derives its first exact intervals.
Assume throughout this section that $p\ge5$, $p\nmid e$, and $\mu_p\subset K$. After completed unramified base change, take the common tame good field $L/K$ of degree $N=p^2-1$ and put $s=(p+1)e$.

\begin{lemma}\label{lem:prime-index-irreducible}
A $p^2$-torsion comparison with irreducible residual inertia determines the finite special characters when $p\nmid e$.
\end{lemma}
\begin{proof}
The normal residual wild subgroup is trivial by irreducibility. Over $L$ the first-level module is rational, and irreducibility forces a one-slope polygon. Its first-level coordinate value is $e$ and every primitive second-level value is $e/q$, where $q=p^2$. As in Section~\ref{sec:irreducible}, the first-level leading spaces are scalar multiples of $\F_q$. If the special characters differed, their identification would be Frobenius-semilinear after scalar normalization.

Choose corresponding primitive second-level points. Their common point field $E/L$ is Galois with group a subgroup of the rational first-level translation group. Its degree is at most $q$, and the denominator of $e/q$ forces degree exactly $q$. Thus every first-level translation occurs. In $E$ the two coordinates have value $e$, prime to $p$, and every nonidentity displacement has value $qe$. With a fixed uniformizer and fixed coefficient representatives, write the leading coordinates as $a_i\varpi^e$. The exact prime-valuation formula gives common break $Ne$ and leading displacement $\bar e a_i d_g$ for every $g\ne1$. Consequently the two first-level leading maps on the entire translation group are scalar multiples of the same map $g\mapsto d_g$.

Their prescribed Frobenius-semilinear identification would therefore make Frobenius a fixed scalar multiplication on a scalar multiple of all $\F_q$. This is impossible: after rescaling it would assert $z^p=c z$ for all $z\in\F_q$, although that nonzero polynomial has degree $p<q$. Hence the special characters agree. The argument uses a degree-$q$ point field, independently of the degree of the full second division field.
\end{proof}

\begin{lemma}\label{lem:rational-high-recovery}
For two two-slope groups with first-level torsion rational over $L$, every $p^2$-torsion isomorphism identifies their high lines. There is no upper bound on $e$ in this assertion.
\end{lemma}
\begin{proof}
Write their parameters as $D_i$. Rationality gives $p\mid D_i$, so $D_i\ne e$. Suppose the high lines are not identified.

If $D_1<e$, choose a high inverse of a nonzero high point in the first model. Its value is $(s-pD_1)/p$. The $p$ roots of this value form one coset of the high line: translation by a high point preserves the value, whereas translation by a low point lowers it to $D_1/p$. The numerator $s-pD_1$ is prime to $p$, so the point field has degree $p$ and translation group exactly the high line. Its coordinate has normalized value $s-pD_1$, and its nonzero displacements have value $p(s-D_1)$. Its exact break is therefore $eN$.

The corresponding second-model point has a low target and value $D_2/p^3$ over $L$. In the same degree-$p$ field its value is $D_2/q$; nonintegrality is already impossible. Otherwise every nonidentity translation is low in the second model and has displacement $D_2$. The coordinate inequality would give
\[
                  D_2\ge D_2/q+eN,
                  \qquad D_2\ge qe,
\]
contrary to $D_2<pe$.

It remains that both $D_i>e$; reorder so $D_1\le D_2$. A high-target inverse in the first model has value $(s-D_1)/q$, with prime-to-$p$ numerator. Its common field has degree $q$ and every first-level translation occurs. The two normalized values are $s-D_1$ and $D_2/p$. A translation high in the first model and low in the second has exact first-model break $N(s-D_1)$. The second-model inequality implies $D_2/p\ge s-D_1$. But
\[
                    s-D_1\ge s-D_2>D_2/p
\]
because $D_2<pe$. This is the final contradiction. No choice of a complementary tame eigenline is involved, so scalar residual action causes no exception.
\end{proof}

\begin{lemma}\label{lem:one-slope-recovery}
A $p^3$-torsion comparison cannot have one one-slope polygon and one two-slope polygon. If both polygons have one slope, their special characters agree.
\end{lemma}
\begin{proof}
For a one-slope group the first-level leading map, obtained by dividing a coordinate by $\pi_L^e$, is injective and linear. Wild inertia over $L$ fixes it, so first-level torsion is $L$-rational. Tame inertia acts on its two-dimensional image by scalar multiplication by $\lambda=\zeta^{c+e}$. If the residual representation is reducible, an invariant line makes $\lambda\in\F_p$, and the entire residual action is scalar. If both polygons are one-slope, the irreducible case is Lemma~\ref{lem:prime-index-irreducible}; in the reducible case equality of these actual scalars gives equal $c$.

Suppose one polygon is two-slope with parameter $B$. Its high line forces the common residual representation to be reducible, hence scalar. In particular $p\mid B$. Transport that high line to the one-slope group and take both actual quotients; a $p^2$-comparison remains. The one-slope quotient has parameter $e$, prime to $p$. The other quotient has parameter $pB$ if $B<e$, or $s-B$ if $B>e$. In the first case opposite divisibility contradicts Lemma~\ref{lem:prime-residual-comparison}. In the second both parameters are prime, and the first-level break formula forces $s-B=e$, hence $B=pe$, contrary to the strict two-slope range. The equality $B=e$ is excluded by divisibility. These cases exhaust the possibilities.
\end{proof}

\begin{theorem}\label{thm:prime-index-upper}
Let $n\ge3$ and $e<p^n$. In the full height-two filtration class over a complete field with perfect residue field, $p\ge5$, $p\nmid e$, and $\mu_p\subset K$, a $p^n$-torsion comparison determines the finite special character. Equivalently,
\[
                         \nuMin(p,e,2)\le n-1.
\]
\end{theorem}
\begin{proof}
We first compare two whole supersingular groups. The one-slope cases and mixed polygons are covered by Lemma~\ref{lem:one-slope-recovery}. For two two-slope groups, nontrivial residual wild inertia identifies the high line as its unique fixed line. If the residual action is tame, its stable high line has a stable complement, both characters take values in $\F_p^\times$, and $L$ kills them. Thus their first-level modules are rational, and Lemma~\ref{lem:rational-high-recovery} identifies the high lines from the available $p^2$-comparison.

Write the original parameters and tangent exponents as $B_i,c_i$. Equality of the determinant and high-line characters gives
\begin{equation}\label{eq:prime-index-congruences}
 c_1\equiv c_2\pmod{p-1},\qquad
 c_1-B_1\equiv c_2-B_2\pmod N.
\end{equation}
If at least one parameter is prime to $p$, the residual action is wild and Lemma~\ref{lem:prime-residual-comparison} recovers the character. Assume both are divisible by $p$. They cannot equal $e$. If both exceed $e$, their high quotients have parameters $s-B_i$ prime to $p$. Their first-level comparison suffices: a low point field has degree $p$ and exact break $ps-(p+1)(s-B_i)$, so the parameters and high characters agree. If the original parameters straddle $e$, the next first-level modules have trivial and nontrivial wild action respectively. Indeed the small high quotient has two distinct stable lines, its old image and its high line, whereas the larger quotient has nonintegral low value $(s-B_i)/p$. This is impossible.

The remaining case is $B_i=pb_i<e$. We iterate actual high quotients. After $j$ quotients the available level is exactly $p^{n-j}$ and, while all preceding parameters are below $e$, the current parameter is
\[
                         D_i^{(j)}=p^{j+1}b_i.
\]
At every stage $1\le j\le n-2$, the preceding small quotient supplies a stable low line distinct from the current high line. Thus the current first-level module is $L$-rational. Its available level is at least $p^2$, so Lemma~\ref{lem:rational-high-recovery} identifies the high lines and permits the next quotient. This next comparison retains at least $p$-torsion, even at $j=n-2$.

At any such stage the two divisible parameters cannot equal $e$. If both exceed $e$, the next prime parameters are recovered from their first-level breaks. If they straddle $e$, the next residual wild actions disagree. The only remaining possibility is that both stay below $e$ through $j=n-2$. Then
\[
                 0<b_i<e/p^{n-1}<p,
                 \qquad b_i\in\{1,\ldots,p-1\}.
\]
Equation~\eqref{eq:prime-index-congruences} gives $b_1\equiv b_2\pmod{p-1}$. Reduction modulo $p-1$ is injective on that precise interval, so $b_1=b_2$ and the common original high character gives $c_1=c_2$ modulo $N$. This completes the whole supersingular comparison, with every loss of a torsion level accounted for.

For the full filtration class, pass to the completed maximal unramified base. Since $p\nmid e$, the cyclotomic character modulo $p^n$ has wild order $p^{n-1}$. The determinant identifies total dimensions and finite determinant factors. Dimensions zero and two, and two groups filtered by height-one pieces, are determined residually as before. A mixed filtered/supersingular comparison would require $e>p^{2n-3}$ by Proposition~\ref{prop:mixed-obstruction}, whereas $e<p^n\le p^{2n-3}$ for $n\ge3$. Thus it is excluded. The whole-group argument just proved settles the remaining case.
\end{proof}

\begin{corollary}\label{cor:prime-index-intervals}
Assume $p\ge5$, $p\nmid e$, and $\mu_p\subset K$. Then
\[
\begin{array}{c|c}
\text{absolute ramification interval}&\nuMin(p,e,2)\\\hline
 p<e<p^3&2\\
 p^3<e<p^4&3.
\end{array}
\]
More generally, if $e>p$ and $p^k<e<p^{k+1}$, then
\[
 \left\lfloor\frac{k+3}{2}\right\rfloor
       \le\nuMin(p,e,2)\le\max\{2,k\}.
\]
\end{corollary}
\begin{proof}
Take $n=3$ or $4$ in Theorem~\ref{thm:prime-index-upper}. Theorem~\ref{thm:general-arc} gives failures at $p^2$ when $e>p$, and at $p^3$ when $e>p^3$, proving both exact values. For the general upper bound use $n=\max\{3,k+1\}$. Put $r=\lfloor(k+3)/2\rfloor$; then $2r-3\le k$ and the same construction gives a failure at $p^r$, hence $\nuMin\ge r$. No equality is asserted when these bounds differ.
\end{proof}
